% ---- pg-430 （书页 418） ----

\begin{mvtable}
{\raggedleft\small 续表\par}
\small
\begin{tabular}{l|ccccccc}
\toprule
& \multicolumn{3}{c|}{指\quad 标} & \multicolumn{4}{c}{得\quad 分}\\
\cmidrule(lr){2-4}\cmidrule(l){5-8}
& 销售增长 & 销售利润 & \makecell{新客户\\销售额} & 创造力 & 机械推理 & 抽象推理 & 数学能力\\
& $(x_1)$ & $(x_2)$ & $(x_3)$ & $(x_4)$ & $(x_5)$ & $(x_6)$ & $(x_7)$\\
\midrule
39 & 108.3 & 122.3 & 108.5 & 15 & 19 & 12 & 41\\
40 & 106.8 & 119.0 & 106.8 & 14 & 20 & 12 & 37\\
41 & 102.5 & 109.3 & 103.8 & 09 & 17 & 13 & 32\\
42 & 92.5 & 102.5 & 99.3 & 13 & 15 & 06 & 23\\
43 & 102.8 & 113.8 & 106.8 & 17 & 20 & 10 & 32\\
44 & 83.3 & 87.3 & 96.3 & 01 & 05 & 09 & 15\\
45 & 94.8 & 101.8 & 99.8 & 07 & 16 & 11 & 24\\
46 & 103.5 & 112.0 & 110.8 & 18 & 13 & 12 & 37\\
47 & 89.5 & 96.0 & 97.3 & 07 & 15 & 11 & 14\\
48 & 84.3 & 89.8 & 94.3 & 08 & 08 & 08 & 09\\
49 & 104.3 & 109.5 & 106.5 & 14 & 12 & 12 & 36\\
50 & 106.0 & 118.5 & 105.0 & 12 & 16 & 11 & 39\\
\bottomrule
\end{tabular}
\end{mvtable}

\noindent 9.27\quad 把 $\mathbf{R}$ 因子化以代替样本协方差矩阵, 重复练习 9.26. 并且, 对有标准化体重$[0.8,-0.2,-0.6,1.5]$的老鼠, 用载荷的极大似然估计和\hyperref[eq:9-58]{式(9-58)}计算因子得分.

\noindent 9.28\quad 对\hyperref[tab:1.9]{表 1.9} 的国家和地区女子田径纪录作因子分析. 用样本协方差矩阵 $\mathbf{S}$ 并解释这些因子. 计算因子得分, 并查找这些数据的离群值. 用样本相关系数矩阵 $\mathbf{R}$ 再作一次因子分析. 若把 $\mathbf{R}$ 而非 $\mathbf{S}$ 因子化, 有什么区别吗? 为什么?

\noindent 9.29\quad 参看练习 9.28, 将国家和地区女子田径纪录换为以米/秒为单位的速度测量值(见练习 8.19). 对速度数据作因子分析. 用样本协方差矩阵 $\mathbf{S}$ 并解释这些因子. 计算因子得分, 并查找数据的离群值. 用样本相关系数矩阵 $\mathbf{R}$ 再作一次因子分析. 当 $\mathbf{R}$ 而非 $\mathbf{S}$ 因子化时, 结果有不同吗? 请解释. 和练习 9.28 的结果比较一下, 你认为哪一个分析更好? 为什么?

\noindent 9.30\quad 对\hyperref[tab:8.6]{表 8.6} 中的国家和地区男子田径纪录作因子分析, 重复练习 9.28 的步骤. 男子数据的恰当因子模型同女子的有区别吗? 若没有区别, 因子解释是否大致相同? 若认为这两个模型确有区别, 试说明它们的不同之处.

\noindent 9.31\quad 参见练习 9.30. 将国家和地区男子田径纪录换为以米/秒为单位的速度测量值(见练习 8.21). 对速度数据作因子分析. 用样本协方差矩阵 $\mathbf{S}$ 并解释这些因子. 计算因子得分, 并查找数据的离群值. 用样本相关系数矩阵 $\mathbf{R}$ 再作因子分析. 当 $\mathbf{R}$ 而非 $\mathbf{S}$ 因子化, 结果有不同吗? 请说明. 与练习 9.30 的结果比较一下, 你认为哪一个分析更好? 为什么?

\noindent 9.32\quad 对\hyperref[tab:1.10]{表 1.10} 的公牛的数据作因子分析, 用这 7 个变量 YrHgt, FtFrBody, PrctFFB, Frame, BkFat, SaleHt 和 SaleWt. 因子化样本协方差矩阵 $\mathbf{S}$ 并解释这些因子. 计算因子得分, 并查找离群值. 用样本相关系数矩阵 $\mathbf{R}$ 再作因子分析. 比较分别由 $\mathbf{S}$ 和 $\mathbf{R}$ 得到的结果. 若 $\mathbf{R}$ 而非 $\mathbf{S}$ 因子化, 结果有不同吗? 请说明.

\noindent 9.33\quad 对\hyperref[tab:4.6]{表 4.6} 中的心理轮廓数据, 用 5 个变量 Indep, Supp, Benev, Confirm 和 Leader 的度量值的相关矩阵 $\mathbf{R}$ 作因子分析. 对 $m=2$ 和 $m=3$ 分别用主成分和极大似然法求解. 你能解释因子吗? 分析中要包含因子旋转和因子得分计算.

注: 需知极大似然解可能导致海伍德(Heywood)现象.
